% chapters/17/latihan.tex
\bagian{Latihan dan Refleksi}

\subsection*{Latihan Peta Kompetensi}

\begin{enumerate}
  \item Jelaskan perbedaan pertanyaan yang dijawab oleh CPL, CPMK, Sub-CPMK,
        dan indikator.
  \item Sebuah CPMK memperoleh bobot 22 persen, sedangkan CPL yang menaunginya
        memperoleh 43 persen. Jelaskan bagaimana kedua angka itu diperoleh.
  \item Sebuah mahasiswa memperoleh nilai praktikum 90, kuis 60, UTS 52,
        Mini Proyek 70, Proyek Akhir 80, dan partisipasi 80. Hitunglah nilai
        akhirnya sampai satu angka di belakang koma, lalu tentukan hurufnya.
  \item Berdasarkan nilai pada soal nomor tiga, hitunglah nilai \CPMK{1} dan
        \CPMK{3}. Sub-CPMK mana yang paling perlu diperiksa lebih dahulu, dan
        mengapa?
  \item Jelaskan mengapa Proyek Akhir memperoleh bobot terbesar, walaupun
        pengerjaannya lebih singkat daripada seluruh praktikum.
  \item Sebuah kelompok menuntaskan setiap praktikum dengan nilai tinggi,
        tetapi prototipe akhirnya tidak dapat dimainkan sampai selesai.
        Sebutkan tiga kemungkinan penyebabnya, dan jelaskan cara memeriksa
        masing-masing.
  \item Pemain menembus lantai pada level Unity. Dengan berpedoman pada
        Gambar~\ref{fig:lapisan-permainan}, jelaskan urutan lapisan yang perlu
        diperiksa, dan mengapa urutannya demikian.
  \item Jelaskan mengapa refleksi yang menyatakan ``Unity lebih mudah''
        belum menjawab pertanyaan refleksi pada RPS.
\end{enumerate}

\begin{refleksibox}
Lembar refleksi akhir semester pada Bagian~\ref{sec:refleksi-akhir} memuat
tiga pertanyaan pokok. Jawablah ketiganya di sini, dengan berpedoman pada
Tabel~\ref{tab:refleksi-bukti}: setiap pernyataan sebaiknya memuat keadaan
yang konkret, tindakan yang diambil, dan hasil yang dapat diperiksa.

\begin{enumerate}
  \item \textbf{Konsep tersulit di blok Python/Pygame.} Sebutkan konsepnya,
        gejala salah pahamnya pada pekerjaan Anda, dan usaha yang Anda
        lakukan.

  \item \textbf{Konsep tersulit di blok Unity.} Sebutkan konsepnya, gejala
        salah pahamnya pada pekerjaan Anda, dan usaha yang Anda lakukan.

  \item \textbf{Perbedaan Pygame dan Unity} ditinjau dari kemudahan, kendali
        atas algoritma, dan kecepatan pengembangan. Jawablah ketiga tinjauan
        itu secara terpisah, dengan contoh dari pekerjaan Anda sendiri.

  \item \textbf{Kompetensi yang ingin dikembangkan} setelah mata kuliah ini,
        beserta langkah pertama yang akan Anda mulai dalam sebulan ke depan.
\end{enumerate}
\end{refleksibox}
